\subsection{Querying before the institution is fixed}
\label{sec:r19-query-complexity}
The two-economy separation extends to an exact information requirement for a menu of institutions. Suppose the purchaser is told only that the unknown compact convex response set lies in $K$, where $K$ contains the Euclidean ball $\bar z+R\mathbb B^k$, $R>0$. Private payoff at the center is zero for every response, and a query returns $h_Z(t)=\max_{z\in Z}t\cdot z$. The class includes all nonempty compact convex subsets of that ball. It can be implemented by a finite-horizon contract with a compact first-date action set and absorbing continuation. Unlike the earlier simplex construction, this larger class need not have finitely many actions.

\begin{theorem}[Menu information and institutional flexibility]
\label{thm:r19-query-complexity}
Let $k\geq2$.
\begin{enumerate}
\item Suppose there are $m$ nonzero institutional payoff vectors $w^1,\ldots,w^m$, no two of which are positive multiples. To identify all robust returns $\min_Z w^j\cdot z$ exactly for every admitted response economy, the worst-case number of exact value queries is $m$. One query on each ray $\{-hw^j:h>0\}$ suffices. No deterministic adaptive procedure using fewer queries can provide such universal identification. This is equivalently a requirement for ranking each institution against \emph{every} safe-offer level; it need not hold for a single preassigned safe offer whose rank is already evident.
\item More quantitatively, fix $0<r<R$ and any finite query transcript generated by $Z_r=\bar z+r\mathbb B^k$. For a unit direction $u$ let $q_1,\ldots,q_n$ be its normalized nonzero query directions and $a(u)=\max_i q_i\cdot u$, taking $a(u)=-1$ if there are no such queries. Every $\delta>0$ satisfying
\begin{equation}
 \delta\leq R-r,\qquad (r+\delta)a(u)\leq r
 \label{eq:r19-undetected-cap}
\end{equation}
gives another consistent economy whose robust return in direction $w=-u$ is lower by $\delta$. Against a safe offer halfway between the two returns, the conditional minimax regret is $\delta/4$, even with randomized final selection.
\item Suppose queries are acquired before the institutional direction and safe-offer level are announced. To guarantee regret at most $\varepsilon$ for every subsequent unit direction, every safe level, and every admitted economy, a deterministic adaptive acquisition procedure must, on the ball transcript, choose directions whose spherical caps of angular radius
\begin{equation}
 \alpha=\arccos\frac{r}{r+4\varepsilon}
 \label{eq:r19-angular-radius}
\end{equation}
cover the unit sphere, whenever $0<4\varepsilon<R-r$. In particular,
\begin{equation}
 n\geq\frac{\displaystyle\int_0^\pi\sin^{k-2}\theta\,d\theta}
 {\displaystyle\int_0^\alpha\sin^{k-2}\theta\,d\theta}.
 \label{eq:r19-query-cover}
\end{equation}
For fixed $k$, this necessary query count grows at least on the order of $(r/\varepsilon)^{(k-1)/2}$ as the tolerated regret tends to zero.
\end{enumerate}
The exact-query upper bound in the first part applies after the institution menu is known. The last part instead prices the flexibility to choose an arbitrary institutional direction after information acquisition. It is an information bound, not a lower bound on the computational cost of solving a known dynamic model.
\end{theorem}

The proof is in the supplementary appendix. The argument does not treat arbitrary real-valued observations as having a fixed bit budget. Rather, it exploits the restricted geometry of a value query: an additional response can change a purchaser's adverse return while remaining hidden behind every queried supporting plane. In the first part, knowing the relevant institutional rays removes this ambiguity with exactly one message per ray. In the last part, information must also be prepared for directions not yet chosen. A purchaser who changes fee incidence or quality penalties after acquiring an oracle may therefore need genuinely new dynamic value information, not simply a more accurate version of an old coordinate message. The sphere bound allows signed payoff directions, including penalties as well as rewards; restricting the eventual institutional directions restricts the region that must be covered.
